% !TeX root = ../navier-stokes-manuscript.tex
% ============================================================
% 1.9 One Fluid, One Coherent Field, and the VPI Law
% ============================================================
\subsection{One Fluid, One Coherent Field, and the VPI Law}\label{sub:OneFluidOneCoherentFieldAndTheVPILaw}

Dead, Jump, and Blown all ended at the same physical break: some portion of the fluid no longer answers the information coming from its neighbours.
To know exactly what has broken, we have to say what that portion was meant to remain inside.
It was never an independent point.
It was part of one velocity field, carried by transport while viscosity, pressure, and incompressibility joined its next motion to the motion around it.
Following a material packet makes that relation literal.

Take a small parcel of fluid at time zero and let the flow map \(X(a,t)\) carry each initial material label \(a\):
\[
  \partial_tX(a,t)=u(X(a,t),t),
  \qquad
  X(a,0)=a.
\]
The velocity field moves the parcel and simultaneously changes the distances, directions, and shapes in its material neighbourhood.
The matter remains part of the same fluid while its geometry is continually rearranged.

Incompressibility fixes what this rearrangement can do to volume.
If
\[
  \mathcal V(a,t):=\det\nabla_aX(a,t),
\]
then
\[
  \frac{d}{dt}\mathcal V(a,t)
  =
  (\nabla\cdot u)(X(a,t),t)\mathcal V(a,t).
\]
Since \(\nabla\cdot u=0\),
\[
  \mathcal V(a,t)=\mathcal V(a,0)=1.
\]
The packet may stretch, turn, and fold through space while its material volume remains unchanged.

Its carried velocity is
\[
  U(a,t):=u(X(a,t),t),
\]
and its acceleration is
\[
  \frac{d}{dt}U(a,t)
  =
  -\nabla p(X(a,t),t)
  +
  \nu\Delta u(X(a,t),t).
\]
Transport has already chosen the point and neighbourhood at which pressure and viscosity act.
Both forces determine the next velocity of the same parcel at the same instant.

The Laplacian gives viscosity a direct local meaning.
For an interior point of a smooth three-dimensional field, the mean-value expansion over a ball gives component by component
\[
  \Delta u(x,t)
  =
  \lim_{r\downarrow0}
  \frac{10}{r^2}
  \left(
    \frac{1}{|B_r(x)|}
    \int_{B_r(x)}u(y,t)\,dy
    -
    u(x,t)
  \right).
\]
Viscosity reads how the velocity at \(x\) curves away from the average velocities in smaller and smaller spatial neighbourhoods.
As transport changes those neighbourhoods, the local comparison acting on the parcel changes with them.

Pressure reaches through the fluid in a different way.
Taking the divergence of momentum gives
\[
  -\Delta p
  =
  \partial_i\partial_j(u_i u_j).
\]
For the decaying whole-space field, the decay normalization fixes pressure and gives
\[
  p
  =
  R_iR_j(u_i u_j).
\]
Changing velocity far from the parcel can change the pressure gradient at the parcel, because pressure enforces the divergence-free constraint across the whole field.

These are simultaneous parts of the VPI law.
VPI means viscosity, pressure, and incompressibility acting through one transported velocity:
\[
  \begin{gathered}
  \partial_tX=u(X,t),\\
  \frac{d}{dt}u(X,t)
  =
  -\nabla p(X,t)
  +
  \nu\Delta u(X,t),\\
  \nabla\cdot u=0,
  \qquad
  -\Delta p=\partial_i\partial_j(u_i u_j).
  \end{gathered}
\]
Viscosity couples the parcel to the curvature of its immediate velocity neighbourhood, pressure couples it to the configuration of the whole fluid, and incompressibility constrains how every material neighbourhood can move.
Transport changes the geometry on which all three act and carries their resulting velocity into the next instant.
Together they form one law of one fluid, with every mechanism acting on the others in the same instant.

The same statement holds for a packet with a moving boundary.
Let \(A(t)=X_t(A(0))\) be a smooth material region.
The Reynolds transport theorem and the momentum equation give
\[
  \frac{d}{dt}\int_{A(t)}u(x,t)\,dx
  =
  \int_{\partial A(t)}
  \left[
    -pI
    +
    \nu\bigl(\nabla u+\nabla u^{\mathsf T}\bigr)
  \right]n\,dS.
\]
Pressure traction and viscous traction act together across the packet boundary, while incompressibility preserves the volume being carried.
The packet momentum changes through the action of the surrounding part of the same fluid.

Vorticity and strain give another reading of this response without introducing another object.
Write
\[
  \nabla u=S+\Omega,
  \qquad
  S:=\frac12\bigl(\nabla u+\nabla u^{\mathsf T}\bigr),
  \qquad
  \omega:=\nabla\times u.
\]
For an incompressible smooth solution,
\[
  D_t\omega
  =
  S\omega+\nu\Delta\omega.
\]
The strain generated by the velocity gradient stretches and turns the vorticity carried by the flow, while viscosity compares that vorticity with its local surroundings.
Pressure is still determined by the velocity field through the normalized Poisson relation, and the flow map still preserves volume.
Velocity, vorticity, strain, pressure, and the derivative tower are different readings of the same coherent field.

\divider{A Moving Material Graph}

The packet relation can be pictured, but not replaced, by a moving graph.
Carry a shape-regular family of small material cells \(A_a^{(h)}(t)\) by the exact flow map, use each cell as a vertex, and join cells that share a material face.
Let \(U_a^{(h)}\) and \(X_a^{(h)}\) be the average velocity and centroid of cell \(a\), and let \(\tau_{ab}^{(h)}(t)\) be the transmissibilities of a consistent moving finite-volume discretization.
Then, as the mesh size \(h\) tends to zero,
\[
  \bigl(L_t^{(h)}U^{(h)}\bigr)_a
  :=
  \frac{1}{|A_a^{(h)}(t)|}
  \sum_{b\sim a}
  \tau_{ab}^{(h)}(t)
  \bigl(U_b^{(h)}(t)-U_a^{(h)}(t)\bigr)
  =
  \Delta u\bigl(X_a^{(h)}(t),t\bigr)+o(1).
\]
The transmissibilities move with the packet geometry, pressure at one vertex still depends on the velocity configuration across all the cells, and incompressibility still constrains every cell volume.
This is only a discrete picture of the continuum VPI law above.
No vertex chooses its response independently, and the graph adds no second fluid object.

At every finite rung of the derivative tower, transport, pressure, viscosity, and incompressibility remain equations of this same field.
Dead, Jump, and Blown describe ways a portion of that field can cease to join the response around it, while the fluid on every side of the failure remains the same one fluid.
A shrinking structure can carry less kinetic energy while its derivative responses intensify, so the one-fluid picture still needs a measurement that remains visible under the change of scale and stays tied to the finite time available.
